% authority: science_draft
\documentclass[11pt]{article}

\usepackage[margin=1in]{geometry}
\usepackage{amsmath,amssymb,amsthm}
\usepackage{booktabs,longtable,array,enumitem}

\newtheorem{definition}{Definition}
\newtheorem{theorem}{Theorem}
\newtheorem{proposition}{Proposition}
\newtheorem{obstruction}{Precise obstruction}
\newenvironment{nonconclusion}{\par\smallskip\noindent\textbf{Non-conclusion.}\ }{\par\smallskip}

\newcommand{\Q}{\mathbb Q}
\newcommand{\Vsrc}{\mathcal V_{\mathrm{src}}}
\newcommand{\Dsrc}{\mathcal D_{\mathrm{src}}}
\newcommand{\Jsrc}{\mathcal J_{\mathrm{src}}}
\newcommand{\One}{\mathbf 1}

\title{SourceKernelDirichletVariationalCandidate\_v1\\
\large A finite source variational object, exact flux balance, and a scoped
directional-update obstruction}
\author{Ontology Formalizer Packet RT-20260728-005}
\date{2026-07-28}

\begin{document}
\maketitle

\section{Status and decisive result}

This task-local object is \texttt{draft/control}, \texttt{proposal-only},
and \texttt{source-extension data}.  It executes v21 P7-T06 under
\texttt{CB-V21-P7-T06-SOURCE-VARIATIONAL-OBJECT-001}.  It does not modify
canonical ontology, adopt a source law, or promote a physical interpretation.

For every declared finite rational P7-T02 row kernel \(P\), the packet
constructs the symmetric conductance
\[
 C=\frac12(P+P^{\mathsf T})
\]
and the dimensionless source functional
\[
 \Vsrc(q)=\frac14\sum_{x,y}C_{xy}(q_y-q_x)^2 .
\]
Its first variation is exactly a symmetric graph-Laplacian operator.  The
associated antisymmetric edge record has zero total divergence on a closed
finite carrier and an exact boundary-flux formula on every subset.

The same construction deliberately discards the directional residue
\[
 A=\frac12(P-P^{\mathsf T}).
\]
For the fixed P7-T02 forward control,
\[
 P_{\mathrm{fwd}}=
 \begin{pmatrix}
 1/2&1/2\\
 0&1
 \end{pmatrix},
 \qquad
 I-P_{\mathrm{fwd}}=
 \begin{pmatrix}
 1/2&-1/2\\
 0&0
 \end{pmatrix}.
\]
The latter operator is nonsymmetric, whereas the Hessian of a twice
differentiable scalar functional in the declared Euclidean pairing is
symmetric.  Therefore this full directional update is not the gradient of a
scalar quadratic functional in the fixed class.

The decisive classification is
\[
 \texttt{precise\_obstruction\_with\_constructive\_restricted\_variational\_object}.
\]
Adoption is \texttt{blocked\_adoption\_open\_continuation}.  The obstruction
is not a global no-go theorem: auxiliary variables, different pairings,
nonlocal functionals, nonscalar principles, or a new source-side law remain
outside its scope.

\section{Fixed source basis and conventions}

Let \(\Omega=\{x_1,\ldots,x_n\}\) be one declared finite P7-T02 carrier.
Write the P7-T02 kernel as a row-oriented matrix
\[
 P_{xy}
 =
 \sum_{\substack{T\in E_x\\t(T)=y}}K_{\mathfrak C}(T\mid x),
 \qquad
 P_{xy}\in\Q_{\geq0},
 \qquad
 \sum_yP_{xy}=1.
\]
This is the transpose convention of the target-row/source-column update
matrix used for bookkeeping states.  It changes no source record and makes
the source-to-target direction explicit.

The entries of \(P\) are dimensionless bookkeeping coefficients.  The row
label is not physical time; \(P\) is not a propagator, transition probability,
Hamiltonian, kinetic operator, or empirical response law.  P7-T02 sector
preservation and the P7-T05 equipped-control comparison remain
proposal-only.  Neither supplies a target atlas, target metric, matter
coupling, action, stress-energy tensor, or physical conservation law.

The new probe
\[
 q:\Omega\longrightarrow\Q
\]
is a dimensionless task-local test coordinate.  In particular, this \(q\)
is \emph{not} the inherited P7-T01 charge-chain record \(q(\Psi)\).  It is
also not a matter field, charge density, clock reading, detector outcome,
spacetime coordinate, or observable.  The standard finite pairing
\(\langle u,v\rangle=\sum_xu_xv_x\) is stipulated only for this candidate.

\section{Symmetric source construction}

\begin{definition}[Conductance and directional residue]
For a declared finite row kernel \(P\), define
\begin{equation}
 C=\frac12(P+P^{\mathsf T}),
 \qquad
 A=\frac12(P-P^{\mathsf T}).
 \label{eq:split}
\end{equation}
Then \(C^{\mathsf T}=C\), \(A^{\mathsf T}=-A\), and \(P=C+A\).
Because \(P\) is entrywise nonnegative, \(C\) is entrywise nonnegative.
The word ``conductance'' names a finite symmetric coefficient record only.
\end{definition}

\begin{definition}[Source variational functional]
For \(q\in\Q^\Omega\), define
\begin{equation}
 \Vsrc(q)
 =
 \frac14\sum_{x,y\in\Omega}C_{xy}(q_y-q_x)^2.
 \label{eq:functional}
\end{equation}
Let \(L_C\) be the symmetric finite Laplacian
\begin{equation}
 (L_C)_{xy}
 =
 \begin{cases}
   -C_{xy}, & x\neq y,\\[2pt]
   \displaystyle\sum_{z\neq x}C_{xz}, & x=y.
 \end{cases}
 \label{eq:laplacian}
\end{equation}
Then \(\Vsrc(q)=\tfrac12 q^{\mathsf T}L_Cq\).  Diagonal entries of \(C\)
cancel from Equation~\eqref{eq:functional}.
\end{definition}

The functional is invariant under the formal shift
\(q\mapsto q+c\One\), is nonnegative over \(\Q\), and vanishes exactly when
\(q\) is constant on every connected component of the undirected support of
\(C\).  These are finite algebraic facts, not physical gauge invariance,
vacuum stability, or energy positivity.

The v21 phrase ``source-equivalent'' is used here only for this source
variational surrogate, which is derived from symmetric kernel coefficients.
It does not mean information equivalence, dynamical equivalence to the P7-T02
update, or equivalence to a physical matter action.  Indeed, if
\(D_C=\operatorname{Diag}(C\One)\), then
\begin{equation}
 (I-P)-L_C=I-D_C-A,
 \qquad
 \operatorname{Sym}(I-P)=I-C.
 \label{eq:information-loss}
\end{equation}
Thus \(L_C=\operatorname{Sym}(I-P)\) only when \(C\One=\One\), and even then
a nonzero \(A\) prevents equality with the full directional operator.  The
functional also ignores every diagonal holding weight because
\(q_x-q_x=0\).

\begin{theorem}[First variation and finite source-flux balance]
\label{thm:variation}
% theorem_id: P7T06-THM-SOURCE-VARIATION-FLUX-BALANCE-001
\begin{center}
\normalfont\small\ttfamily
P7T06-THM-SOURCE-VARIATION-\\
FLUX-BALANCE-001
\end{center}
For every \(q,\eta\in\Q^\Omega\),
\begin{equation}
 \left.\frac{d}{d\epsilon}\right|_{\epsilon=0}
 \Vsrc(q+\epsilon\eta)
 =
 \sum_x\eta_x\Dsrc{}_{x}(q),
 \qquad
 \Dsrc{}_{x}(q)
 =
 \sum_yC_{xy}(q_x-q_y)
 =
 (L_Cq)_x.
 \label{eq:first-variation}
\end{equation}
Define the oriented edge record
\begin{equation}
 \Jsrc{}_{xy}(q)=C_{xy}(q_x-q_y).
 \label{eq:flux}
\end{equation}
Then
\begin{align}
 \Jsrc{}_{xy}&=-\Jsrc{}_{yx},                                  \label{eq:antisym}\\
 \Dsrc{}_{x}&=\sum_y\Jsrc{}_{xy},                              \label{eq:divergence}\\
 \sum_{x\in\Omega}\Dsrc{}_{x}&=0,                              \label{eq:closed}\\
 \sum_{x\in B}\Dsrc{}_{x}
 &=\sum_{\substack{x\in B\\y\notin B}}\Jsrc{}_{xy}
 \quad\text{for every }B\subseteq\Omega.                       \label{eq:boundary}
\end{align}
\end{theorem}

\begin{proof}
Expanding Equation~\eqref{eq:functional} at \(q+\epsilon\eta\) and taking
the coefficient linear in \(\epsilon\) gives
\[
 \frac12\sum_{x,y}C_{xy}(q_y-q_x)(\eta_y-\eta_x).
\]
Exchange \(x\) and \(y\) in the first contribution and use
\(C_{xy}=C_{yx}\).  The result is
\[
 \sum_x\eta_x\sum_yC_{xy}(q_x-q_y),
\]
which proves Equation~\eqref{eq:first-variation}.  Symmetry of \(C\)
immediately gives Equation~\eqref{eq:antisym}, and
Equation~\eqref{eq:divergence} is the definition of \(\Dsrc\).  Summing an
antisymmetric finite array over all ordered pairs proves
Equation~\eqref{eq:closed}.  For a subset \(B\), the terms with both
endpoints in \(B\) cancel in pairs, leaving only the displayed crossing
terms in Equation~\eqref{eq:boundary}.
\end{proof}

\begin{nonconclusion}
\(\Vsrc\) is a source-equivalent variational object only.
\(\Dsrc\) is its finite variational derivative, not a stress-energy tensor.
\(\Jsrc\) is an antisymmetric source-flux record, not a physical current.
Equations~\eqref{eq:closed} and \eqref{eq:boundary} are finite cancellation
identities, not energy-momentum conservation, a Noether theorem,
diffeomorphism invariance, or equations of motion.
\end{nonconclusion}

\section{Boundary, anomaly, and malformed branches}

The closed-carrier identity \(\sum_x\Dsrc{}_{x}=0\) requires only the finite
carrier and symmetric \(C\).  For a retained subset \(B\), its exact
``boundary term'' is the crossing sum in Equation~\eqref{eq:boundary}.
No continuum boundary, normal vector, integration measure, or physical
flux is inferred.

If a proposed source extension adds a declared vector \(a\in\Q^\Omega\),
the task-local residual equation
\[
 \Dsrc(q)=a
\]
has the necessary compatibility condition
\[
 \sum_xa_x=0.
\]
A nonzero total is classified as an \emph{algebraic balance anomaly}
relative to this finite candidate.  It is not a quantum anomaly, dissipative
term, external force, or empirical source.

The construction fails closed when the carrier is infinite, the kernel is
not a rational nonnegative row-normalized P7-T02 record, the probe is
mistyped, the split in Equation~\eqref{eq:split} is not used, or any target
metric, desired GR behavior, detector outcome, validator result, generated
derivative, registry metadata, or checkpoint state is imported as a
mathematical premise.

\section{Directional-update obstruction}

The stronger candidate would require the full directional displacement
\begin{equation}
 F_P(q)=(I-P)q
 \label{eq:directional}
\end{equation}
to equal the Euclidean gradient of one scalar quadratic functional on
\(\Q^\Omega\).

\begin{theorem}[Hessian-symmetry criterion]
\label{thm:hessian}
If \(F(q)=Mq\) is the gradient of a scalar quadratic
\[
 W(q)=\frac12q^{\mathsf T}Hq+b^{\mathsf T}q+c
\]
for all \(q\), using the declared finite Euclidean pairing, then
\(M=H_{\mathrm{sym}}=(H+H^{\mathsf T})/2\) and hence \(M=M^{\mathsf T}\).
Conversely, every symmetric \(M\) is represented by
\(W(q)=\tfrac12q^{\mathsf T}Mq\).
\end{theorem}

\begin{proof}
Direct differentiation gives
\[
 \nabla W(q)=\frac12(H+H^{\mathsf T})q+b.
\]
Equality with the homogeneous map \(Mq\) for all \(q\) forces \(b=0\) and
\(M=(H+H^{\mathsf T})/2\), which is symmetric.  The converse follows by
direct differentiation.
\end{proof}

\begin{obstruction}[Full directional P7-T02 update]
\label{obs:directional}
% theorem_id: P7T06-THM-DIRECTIONAL-UPDATE-VARIATIONAL-OBSTRUCTION-001
% obstruction_id: OBST-P7T06-DIRECTIONAL-KERNEL-NONVARIATIONAL-001
\begin{center}
\normalfont\small\ttfamily
P7T06-THM-DIRECTIONAL-UPDATE-\\
VARIATIONAL-OBSTRUCTION-001\\[2pt]
OBST-P7T06-DIRECTIONAL-KERNEL-\\
NONVARIATIONAL-001
\end{center}
For the fixed positive P7-T02 control on \(\{x_u,x_v\}\),
\[
 P_{\mathrm{fwd}}
 =
 \begin{pmatrix}
 1/2&1/2\\
 0&1
 \end{pmatrix}.
\]
Therefore
\[
 I-P_{\mathrm{fwd}}
 =
 \begin{pmatrix}
 1/2&-1/2\\
 0&0
 \end{pmatrix}
 \neq
 \begin{pmatrix}
 1/2&0\\
 -1/2&0
 \end{pmatrix}
 =
 (I-P_{\mathrm{fwd}})^{\mathsf T}.
\]
By Theorem~\ref{thm:hessian}, no scalar quadratic functional on the declared
two-dimensional probe space with the declared Euclidean pairing has
gradient \((I-P_{\mathrm{fwd}})q\) for every \(q\).
\end{obstruction}

For this control,
\[
 C_{\mathrm{fwd}}
 =
 \begin{pmatrix}
 1/2&1/4\\
 1/4&1
 \end{pmatrix},
 \qquad
 A_{\mathrm{fwd}}
 =
 \begin{pmatrix}
 0&1/4\\
 -1/4&0
 \end{pmatrix},
\]
and
\[
 L_{C_{\mathrm{fwd}}}
 =
 \begin{pmatrix}
 1/4&-1/4\\
 -1/4&1/4
 \end{pmatrix}.
\]
Thus the constructive functional retains an exact symmetric variation while
the nonzero \(A_{\mathrm{fwd}}\) records directional information that it
does not reproduce.  Here \(C_{\mathrm{fwd}}\One=(3/4,5/4)^{\mathsf T}\),
so \(L_{C_{\mathrm{fwd}}}\) is not even
\(\operatorname{Sym}(I-P_{\mathrm{fwd}})\); the degree imbalance in
Equation~\eqref{eq:information-loss} is additional omitted information.

\begin{nonconclusion}
The obstruction is limited to representing the fixed linear operator
\(I-P_{\mathrm{fwd}}\) as the Euclidean gradient of a scalar quadratic on
the same probe variables.  It does not rule out a reversible subfamily,
auxiliary variables, constrained or nonscalar variational principles,
different lawful pairings, nonlocal constructions, a conservative
source-side extension, physical matter actions, or the theory as a whole.
\end{nonconclusion}

\section{Exact status relative to P7-T06}

The construction supplies one reproducible source-equivalent variational
object, one finite variational derivative, one antisymmetric flux record,
one closed-carrier balance identity, one boundary formula, and one algebraic
anomaly condition.  It also supplies a precise obstruction to preserving the
full directional P7-T02 update in the fixed scalar-quadratic class.

It does not satisfy the stronger physical branches of the P7-T06 title.
No reconstructed physical matter variable or source-derived geometry is
available; therefore no effective matter action, spacetime variation,
stress-energy tensor, covariant divergence, physical Noether current, or
Einstein-equation input is derived.  P7 has materially advanced at proposal
scope, while the physical matter-coupling milestone remains open.

The next bounded step is the planned P7-T07 smuggling audit.  It must test
whether the entire P7-T01 through P7-T06 package preserves the source-only
and model-to-world boundaries before any stronger interpretation is
considered.

\section{Authority boundary}

Nothing in this packet authorizes:
\begin{itemize}[leftmargin=*]
\item canonical ontology, source-law, variational-law, or coupling-law
adoption;
\item a physical matter action, stress-energy tensor, conservation law,
interaction, detector semantics, cone, metric, or effective geometry;
\item Einstein equations, exact-GR recovery, benchmark promotion, proof,
publication, push, or completed derivation;
\item a global no-go theorem, theory rejection, or impossibility of future
source extension; or
\item treating child review, validator success, registry state, generated
derivatives, or checkpoint success as mathematical or scientific authority.
\end{itemize}

\end{document}
